% Workshop version of docs/paper/draft.md. Build: tectonic main.tex
\documentclass[10pt,twocolumn]{article}
\usepackage[margin=0.75in]{geometry}
\usepackage[T1]{fontenc}
\usepackage{lmodern}
\usepackage{microtype}
\usepackage{booktabs}
\usepackage{amsmath,amssymb}
\usepackage{url}
\usepackage[hidelinks]{hyperref}
\usepackage{enumitem}
\setlist{nosep,leftmargin=*}

\newcommand{\Sim}{\textsuperscript{S}}
\newcommand{\Cal}{\textsuperscript{C}}
\newcommand{\Meas}{\textsuperscript{M}}

\title{Spending the Spot Notice: Deadline-Bounded Evacuation\\of Agent-Session KV Caches}
\author{Anonymous Author(s)}
\date{}

\begin{document}
\maketitle

\begin{abstract}
Agentic LLM workloads keep tens of thousands of tokens of KV cache per session alive
\emph{between} turns. When a spot GPU node is reclaimed, that cache is lost silently:
no request fails, but every later turn of every resident session pays a full
re-prefill. Existing fault-tolerance systems protect in-flight requests, and agent-KV
systems protect idle sessions only against eviction. Replaying 393 real Claude Code
sessions, we find that 75--78\% of a large serving node's resident KV belongs to
sessions with no turn in flight. We treat the cloud's termination notice as a budget
and plan which session caches to save, and where, under a deadline, bandwidth and
failure domains. In simulation calibrated on an NVIDIA L40S, on-notice evacuation
recovers 47--52\% of reclaim-induced recompute at a 2-minute notice with zero
steady-state cost, within 5--10 points of continuous replication. When replication's
DRAM is instead given to the primary cache (equal memory), evacuation still recovers
about half, while replication does worse than doing nothing. We implement the policy
as an out-of-tree plugin for vLLM's tiered KV offloading. On AWS it evacuates 10.2\,GB
of session KV from an L40S node to a CPU memory node's DRAM (served over NFS) in
6.6\,s, and restores it with greedy output identical to recomputation. For 8B on an
L40S, restore and recompute are at parity (single runs), which we report as a
crossover, not a win.
\end{abstract}

\noindent\emph{Labels:} \Sim{} simulated; \Cal{} simulated with constants measured on
hardware; \Meas{} measured on hardware.

\section{Introduction}
Serving agents is dominated by long, multi-turn sessions. In the AgentX corpus of 393
Claude Code sessions~\cite{agentx}, median input context is about 89k tokens, a
session makes a median of 70 requests, and the median session lasts 1.8\,h.
Providers increasingly serve on preemptible capacity, which AWS reclaims with a
2-minute notice (Azure and GCP give 30\,s).

What does a reclaim cost? Spot and fault-tolerant serving
systems~\cite{spotserve,shuntserve,kevlarflow,lumen,ghostserve} answer for the
\emph{in-flight request}. They restart it, migrate it, or replicate its KV
continuously. For agentic serving this misses the main asset: the \emph{idle session
prefix}. It is not running, so losing it fails nothing, but the next turn silently
re-prefills 50--100k tokens. Agent-KV systems~\cite{continuum,cachescout,kvflow} do
manage idle prefixes, but against eviction on a live node, not node loss.

\paragraph{Contributions.}
(1)~A characterization: on a large node, 75--78\% of resident KV is not running, and
about a third belongs to sessions that never return (\S\ref{sec:char}).
(2)~A deadline-bounded evacuation planner over in-flight and idle KV objects with
release times, per-destination ingress, failure domains and $P(\text{reuse})$. It
generalises ShuntServe's stated future work from requests to
sessions~(\S\ref{sec:design}).
(3)~A regime map of when evacuation matches continuous replication, when it beats it,
and when it cannot help (\S\ref{sec:regime}).
(4)~An out-of-tree vLLM~\cite{vllm} prototype, evaluated on a real L40S, flushing to
local NVMe and to a separate CPU memory node over the network, with output parity
(\S\ref{sec:proto}). The hardware runs exercise the mechanism; the planner's choices
under a tight budget are evaluated only in simulation.

\paragraph{Not claimed.} Predicting preemption~\cite{parcae,skynomad,hazardserve};
exploiting correlated preemption~\cite{skyserve}; that KV transfer generally beats
recompute (ShuntServe~\cite{shuntserve} measured the opposite on g6/g6e, and so do we
on slow tiers); a better reuse predictor~\cite{cachescout,continuum}.

\section{Background}
\begin{table*}[t]
\centering\footnotesize\setlength{\tabcolsep}{3.5pt}
\begin{tabular}{@{}lll@{}}
\toprule
System & Protects & Idle session KV \\
\midrule
SpotServe, ShuntServe & in-flight, at notice & no \\
KevlarFlow, LUMEN, GhostServe & in-flight, continuously & no \\
hazardserve & in-flight, speculatively & no \\
Mooncake, llm-d P2P & prefixes, continuously & lost with holder \\
Continuum, CacheScout & idle, vs.\ eviction & no node loss \\
\textbf{This work} & \textbf{idle + in-flight, at notice} & \textbf{yes} \\
\bottomrule
\end{tabular}
\caption{What each line of work protects, and when.}
\label{tab:related}
\end{table*}
Table~\ref{tab:related} summarises the related work. llm-d's graceful shutdown drains
only in-flight requests. Its P2P sharing~\cite{llmdp2p} is a natural restore path for
evacuated blocks, but it assumes the peer survives. TraceLab~\cite{tracelab} calls for
KV management around the idle gaps in coding-agent workloads.

\section{Design}\label{sec:design}
\paragraph{Budget.} A notice of $W$ seconds and egress bandwidth $B$ against resident
KV $R$ give the budget ratio $\rho = WB/R$. When $\rho \ge 1$ everything fits; when
$\rho < 1$ the planner decides.

\paragraph{Planner.} Objects are in-flight requests (which may finish locally if their
remaining decode fits in $W$) and idle session prefixes. Each has a size $c$, a
recompute cost $r(c) = \frac{c}{p}\left(1 + \frac{c}{2k}\right)$ and a reuse
probability. Its value is $P(\text{reuse})\,(r(c) - \text{reload})$. Destinations
carry a failure domain, free capacity and ingress already booked by other evacuating
nodes. Peers in the reclaim wave are excluded. A greedy pass by value density over one
serial egress pipe with release times places each object on the destination that
finishes it before the deadline at the highest value. Leftover time carries partial
prefixes, since a prefix is usable from the front. Sessions then expand into vLLM
block-hash chains, in prefix order, deduplicated and cut to budget.

\paragraph{vLLM plugin.} vLLM's tiered offloading cascades every stored block to all
secondary tiers, and drops cascades under backpressure. Our offloading spec loads out
of tree, with no fork, and adds: a \emph{deferred} mode, where blocks stay in the CPU
tier and are tracked per agent session; a \emph{hybrid} mode, which writes through
where backpressure admits and tracks the rest as dirty; and a deadline flush. On
notice, a handler writes a flush request \emph{before} pausing admission. The
scheduler-side manager plans and submits store jobs through vLLM's own store path.
Restore needs nothing new: a later request finds the blocks by hash.

\section{Characterization}\label{sec:char}
A discrete-event simulator of a 12-node fleet in 3 failure domains replays the AgentX
corpus, resampled to a load whose median TTFT matches a calibrated baseline. Real
sessions last about 3$\times$ longer than a generator fitted to the same marginals,
so equal session counts overload the fleet. Table~\ref{tab:char} gives steady-state
snapshots.
\begin{table}[t]
\centering\footnotesize\setlength{\tabcolsep}{3.5pt}
\begin{tabular}{@{}lrrr@{}}
\toprule
Node (cache) & KV p50 & Not running & Returning \\
\midrule
8B, 4M tok & 441\,GB & \textbf{78\%} & 66\% \\
70B on 4$\times$H100, 4M tok & 1{,}095\,GB & \textbf{75\%} & 68\% \\
8B, g6e.xlarge, 400k tok & 33\,GB & 100\% & 100\% \\
\bottomrule
\end{tabular}
\caption{Value at risk per node on real sessions\Sim. ``Returning'' is the share of
rebuild cost belonging to sessions that send another turn.}
\label{tab:char}
\end{table}
Three-quarters of what a reclaim destroys is beyond the reach of in-flight fault
tolerance. A third of it is dead weight a planner should skip. And a single
g6e.xlarge has $\rho \ge 1$ in every snapshot: small nodes need no planner.

\section{Regime map}\label{sec:regime}
\begin{table}[t]
\centering\footnotesize\setlength{\tabcolsep}{3.5pt}
\begin{tabular}{@{}lrrrr@{}}
\toprule
Policy & ind.\,30s & ind.\,120s & cor.\,30s & cor.\,120s \\
\midrule
Evacuate (ours) & $-14$ & $\mathbf{-52}$ & $-10$ & $\mathbf{-47}$ \\
Session write-through & $-63$ & $-59$ & $-55$ & $-50$ \\
Hybrid (ours) & $-71$ & $-66$ & $-51$ & $-56$ \\
In-flight replication & $-31$ & $-29$ & $-27$ & $-30$ \\
\bottomrule
\end{tabular}
\caption{Reclaim-induced recompute, \% change vs.\ no reaction, on real sessions with
L40S constants\Cal. Paired seeds; 95\% CIs $\pm$5--40 points. Steady-state egress:
evacuate 0; write-through and hybrid $\sim$50\,TB; in-flight replication
$\sim$584\,TB.}
\label{tab:regime}
\end{table}
Table~\ref{tab:regime} compares policies. Evacuation \emph{wins} in two regimes\Sim.
\emph{At equal memory}, if the DRAM that write-through spends on replicas goes to the
primary cache instead, evacuation reaches $-50\%$ ($-46\%$ correlated) at
$W{=}120$\,s, while write-through \emph{increases} reclaim cost by 84--94\%:
halving the primary cache costs more in steady-state misses than replicas save at
reclaim. This holds only under that memory constraint; with spare DRAM
(Table~\ref{tab:regime}) write-through is the strongest single policy.
\emph{With a shared store} of $\sim$1\,GB/s across 12 nodes, write-through drops
bursts and falls to $-25\%$, while evacuation holds $-43\%$ and hybrid $-51\%$.
Evacuation \emph{cannot help} at $W{=}0$, and recovers only 10--27\% at 30\,s. An
oracle reuse predictor adds no significant benefit over recency on real
sessions\Sim. A synthetic-only finding that failure-domain-blind pools are fragile
did not replicate on real sessions, and we retract it.

\section{Prototype on real GPUs}\label{sec:proto}
\paragraph{Calibration\Meas.} On an L40S with Llama-3.1-8B and vLLM 0.31.0, cold
prefill runs at 16{,}011\,tok/s with a 20.5k-token attention knee (the fit is within
1\% from 32k to 120k tokens), and decode at 46.7\,tok/s. Receiver-side interference
from 1--16\,GB/s of incoming KV traffic is 0--3\%.

\paragraph{Mechanism\Meas.} Two sessions (57k and 20k tokens). Process A runs in
deferred mode, flushes on an emulated 120\,s notice, and is \texttt{SIGKILL}ed. A
fresh process restores turn 2 from the tier, and another computes it cold
(Table~\ref{tab:proto}). Deferred mode left the tier empty until the notice in every
run. Flush projections were within 6--15\%, and restored KV decoded identically to
recomputed KV under greedy decoding. The prompts are random words, so this checks KV
correctness, not answer quality. The memory node stands in for a surviving peer; it is
not a second GPU server, and NFS into vLLM's CPU tier is not a production restore path.

\paragraph{The planner had nothing to decide.} Both sessions fit in 6.6\,s of a 120\,s
notice ($\rho \gg 1$), so every run saved everything. These runs validate the flush
mechanism and its time projection, not the planner's choices, which rest on
\S\ref{sec:regime}.
\begin{table}[t]
\centering\footnotesize\setlength{\tabcolsep}{3.5pt}
\begin{tabular}{@{}lrrc@{}}
\toprule
Tier & Flush (proj.) & TTFT 59k: rest./cold & Par. \\
\midrule
Local NVMe & 49.2 (46.5)\,s & 19.4 / 9.3\,s & \checkmark \\
Memory node & \textbf{6.6} (7.8)\,s & 8.3 / 9.1\,s & \checkmark \\
\bottomrule
\end{tabular}
\caption{Evacuating and restoring 10.17\,GB of session KV\Meas, single runs. Memory node: an r7i.xlarge CPU instance exporting
20\,GB of RAM over NFS in the same zone. At 22k tokens,
restore loses on both tiers (6.6 and 2.9\,s vs.\ 2.4\,s cold).}
\label{tab:proto}
\end{table}
\paragraph{Restore vs.\ recompute.} From local NVMe ($\sim$0.4\,GB/s end to end)
restore loses by 2--3$\times$. From the memory node ($\sim$0.9\,GB/s end to end, limited
by vLLM's promotion pipeline rather than the 1.3--1.5\,GB/s network) it is 9\% faster
at 59k tokens and 23\% slower at 22k. These are single runs. Cold TTFT for the same 59k
prompt was 9.33 and 9.11\,s on two instances, so the 59k gap is suggestive, not
established; we call it parity. This agrees with the arithmetic: for 8B on L40S,
recompute costs the equivalent of $\sim$0.83\,GB/s of KV. We expect larger models to
favour restore, since a 70B model has 2.5$\times$ the KV per token but about
8.75$\times$ the prefill compute, but this is reasoning, not measurement.

\section{Limitations}
The simulator has no memory-based admission control and uses assumed network
constants. Differences under about 20 points are within noise. Spot g6e capacity was
unavailable throughout (placement score 1/10 in three US regions), so the GPU tests
emulate the reclaim with \texttt{SIGKILL}. A real FIS interruption on a CPU spot
instance announced 118\,s of notice, of which at least 97\,s was observed usable
(one run), still about 14$\times$ the 6.6\,s flush. One L4 run forced $\rho \approx
0.3$ with an egress cap and a 30\,s notice: the flush cut mid-session, finished 2\,s
before the deadline (within 0.1\,s of its projection), and a partially evacuated
session restored with output parity. One complete session diverged at its second token; a repeat with logprobs showed an
exact tie there ($-3.0773$ for both candidates), with restored and cold logprobs within
0.01--0.05, so the divergence is floating-point tie-breaking, not corrupted KV.
Its sessions (2k--12k tokens) were below the ${\sim}$54k-token length where restoring
8B KV beats recomputing it, so every restore was valued $\le 0$ and the planner's
value-based choices remain untested on hardware. The GPU timings are single runs without error
bars. GPU results are for 8B only. The ``peer'' was a CPU
memory node serving RAM over NFS; a surviving GPU peer reached by NIXL P2P into GPU
memory would remove the promotion bottleneck, and is unmeasured.

\section{Conclusion}
In agentic serving, most of the state a spot reclaim destroys is idle session KV that
no existing mechanism protects. The termination notice is enough to move it: 10\,GB
in under 7\,s on commodity cloud networking, planned against the deadline, and
restored with output identical to recomputation. Whether moving it \emph{pays} is a
crossover in model size, context length and restore bandwidth. For 8B on an L40S we
measured parity; the larger-model side of the crossover is the next measurement, not a
result.

\bibliographystyle{plain}
\bibliography{refs}
\end{document}
